\documentclass[11pt]{article}
\usepackage[margin=1in]{geometry}
\usepackage{amsmath,amssymb,amsthm}
\usepackage{hyperref}
\usepackage{enumitem}

\newtheorem{theorem}{Theorem}
\newtheorem{lemma}[theorem]{Lemma}
\newtheorem{corollary}[theorem]{Corollary}
\theoremstyle{definition}
\newtheorem{definition}[theorem]{Definition}
\theoremstyle{remark}
\newtheorem{remark}[theorem]{Remark}

\usepackage{xurl}
\let\oldus\_
\renewcommand{\_}{\oldus\allowbreak}

\title{A disproof of Conjecture 00000002291\\[2pt]
\large A centralizer lower bound $|C_G(x)|\ge \tfrac18|G|^{1/2}$ in simple groups}
\date{}

\begin{document}
\sloppy
\maketitle

\section{The conjecture}

Conjecture 00000002291 reads (English text; the Chinese text states the same bound, the same constant
$c=1/8$ and the same example):

\begin{quote}
\emph{Definition:} The Brauer--Fowler theorem bounds centralizers of even-order simple groups:
$|C_G(x)| \le (|G|-1)/2$ type.
\emph{Conjecture:} The exact constant in the lower bound for centralizer sizes:
$C_G(x) \ge c\cdot|G|^{1/2}$ $(c = 1/8)$; $c$ is tight, attained by low-order simple groups
(e.g.\ involutions of $A_5$).
\end{quote}

\paragraph{Reading.} $G$ is a finite simple group of even order, $C_G(x)=\{g\in G: gx=xg\}$ is the
centralizer, $|\cdot|$ is the order and $|G|^{1/2}$ the nonnegative square root. The statement does not
say which $x$ are allowed, so we treat two readings:
\begin{itemize}[leftmargin=2em]
\item[(R1)] $x$ ranges over all nonidentity elements of $G$;
\item[(R2)] $x$ ranges over involutions of $G$ (the Brauer--Fowler context and the $A_5$ example).
\end{itemize}
In both readings we add the hypothesis that $G$ is nonabelian. This only weakens the claim, which makes the
refutation stronger. The second clause, (T), says that $c=1/8$ is tight and attained, with involutions of
$A_5$ named as an example. We do not interpret the preamble's upper-bound remark, and nothing below depends
on it.

\paragraph{What is proved.}
\begin{enumerate}[leftmargin=2em]
\item Under (R1) the lower bound with $c=1/8$ is false: in $A_9$ the $9$-cycle has a centralizer of order at
most $9$, and $9<\tfrac18\sqrt{181440}\approx 53.2$. In fact no constant $c>0$ works under (R1).
\item Clause (T) is false. No element of $A_5$, and so no involution, satisfies $|C(x)|=\tfrac18\sqrt{60}$.
Every nonidentity element of $A_5$ beats the bound by more than a factor of $2$.
\item Under (R2) the lower bound with $c=1/8$ is false: $G=\mathrm{PSL}(2,\mathbb F_{128})$ is simple,
nonabelian and of even order, and the image $t$ of $\begin{pmatrix}1&1\\0&1\end{pmatrix}$ is an involution with
$|C_G(t)|\le128<\tfrac18\sqrt{|G|}\approx181$. In fact no constant $c>0$ works under (R2), using
$\mathrm{PSL}(2,2^k)$ for large $k$.
\end{enumerate}
Since the conjecture is the conjunction of the lower bound and (T), it is false under both readings, and
under each reading Lean refutes the lower bound itself (as well as the named attainment in $A_5$).

\section{Definitions and conventions (as in the Lean file)}

All groups are finite. $A_n=$ \texttt{alternatingGroup (Fin n)}, the even permutations of
$\{0,\dots,n-1\}$; the centralizer is Mathlib's \texttt{Subgroup.centralizer \{x\}}, whose membership is
$g\in C(x)\iff gx=xg$ (\texttt{Subgroup.mem\_centralizer\_singleton\_iff}); orders are \texttt{Nat.card};
the square root is \texttt{Real.sqrt}; an involution is an $x$ with \texttt{orderOf x = 2}. Simplicity is Mathlib's \texttt{IsSimpleGroup} (nontrivial, and the
only normal subgroups are $\bot$ and $\top$).

\begin{definition}[\texttt{LowerBoundAllElements}, reading (R1)]
For every finite group $G$ that is simple, nonabelian and of even order, and every $x\ne 1$ in $G$,
$\tfrac18\sqrt{|G|}\le |C_G(x)|$.
\end{definition}

\begin{definition}[\texttt{LowerBoundInvolutions}, reading (R2)]
For every finite group $G$ that is simple, nonabelian and of even order, and every $x\in G$ of order $2$,
$\tfrac18\sqrt{|G|}\le |C_G(x)|$.
\end{definition}

\begin{definition}[\texttt{AttainedByA5Involution}, clause (T) at $A_5$]
There is $t\in A_5$ with $\operatorname{ord}(t)=2$ and $|C_{A_5}(t)|=\tfrac18\sqrt{|A_5|}$.
\end{definition}

\section{Proofs}

\begin{lemma}[\texttt{finRotate\_mem\_alternatingGroup}, \texttt{card\_centralizer\_finRotate\_perm},
\texttt{card\_centralizer\_cyc\_le}]
Let $n=2m+1\ge 3$ and $\sigma=(0\,1\,\cdots\,n-1)$ (Mathlib's \texttt{finRotate n}). Then $\sigma\in A_n$,
$|C_{S_n}(\sigma)|=n$, and $|C_{A_n}(\sigma)|\le n$.
\end{lemma}
\begin{proof}
$\operatorname{sign}\sigma=(-1)^{n-1}=1$ (\texttt{sign\_finRotate}). The cycle type of $\sigma$ is $\{n\}$
(\texttt{cycleType\_finRotate\_of\_le}). Mathlib's formula \texttt{Equiv.Perm.nat\_card\_centralizer}
gives $|C_{S_n}(g)|=(n-\Sigma)!\cdot\prod\lambda\cdot\prod_k m_k!$ for $g$ of cycle type $\lambda$ with sum
$\Sigma$ and multiplicities $m_k$. For $\lambda=\{n\}$ this is $0!\cdot n\cdot 1!=n$. The map
$C_{A_n}(\sigma)\to C_{S_n}(\sigma)$ forgetting evenness is injective, so $|C_{A_n}(\sigma)|\le n$.
\end{proof}

\begin{lemma}[\texttt{A\_hyps}, \texttt{factorial\_lower}, \texttt{cyc\_ne\_one}]
For $n=2m+1$ with $m\ge 2$: $A_n$ is simple, nonabelian and of even order, $\sigma\ne1$, and
$n!\ge n^2(n-3)$.
\end{lemma}
\begin{proof}
Simplicity is \texttt{alternatingGroup.isSimpleGroup} (needs $n\ge5$). Write $n=(n-3)+3$; then
$n!=n(n-1)(n-2)(n-3)!$, with $(n-1)(n-2)\ge n$ and $(n-3)!\ge n-3$. We have $2|A_n|=n!$
(\texttt{two\_mul\_nat\_card\_alternatingGroup}), and $24=4!\mid n!$, so $|A_n|$ is even. If $A_n$ were
abelian, then $C_{A_n}(\sigma)=A_n$, which would give $n!/2\le n$. That contradicts $n!\ge n^2(n-3)$. If
$\sigma=1$, then $C_{S_n}(\sigma)=S_n$ has order $n!>n$, a contradiction.
\end{proof}

\begin{lemma}[\texttt{ratio\_aux}]
Let $q,r,C,N\in\mathbb N$ and $c>0$ with $q>0$, $C\le q$, $q^2r\le N$ and $c^2r>1$. Then $C<c\sqrt N$.
\end{lemma}
\begin{proof}
$q^2<c^2r\,q^2\le c^2N$, so $C\le q=\sqrt{q^2}<\sqrt{c^2N}=c\sqrt N$.
\end{proof}

\begin{theorem}[\texttt{not\_lowerBoundAllElements}]
Reading (R1) with $c=1/8$ is false.
\end{theorem}
\begin{proof}
Take $G=A_9$ and $x=\sigma$. Then $|G|=9!/2=181440=81\cdot2240$ and $|C_G(x)|\le 9$. Apply
\texttt{ratio\_aux} with $q=9$, $r=2240$, $c=\tfrac18$: $c^2r=35>1$, so $|C_G(x)|<\tfrac18\sqrt{|G|}$.
\end{proof}

\begin{theorem}[\texttt{alternating\_centralizer\_ratio\_small}]
For every real $c>0$ there are $n\ge5$ and $x\ne1$ in $A_n$ such that $A_n$ is simple, nonabelian and of even
order, and $|C_{A_n}(x)|<c\sqrt{|A_n|}$.
\end{theorem}
\begin{proof}
Let $m=\lceil 1/c^2\rceil+2$, $n=2m+1$ and $x=\sigma$. Since $2|A_n|=n!\ge n^2(2m-2)$, we get
$|A_n|\ge n^2(m-1)$, and $c^2(m-1)=c^2(\lceil 1/c^2\rceil+1)>1$. Apply \texttt{ratio\_aux} with
$q=n$, $r=m-1$.
\end{proof}

\begin{theorem}[\texttt{not\_attainedByA5Involution}, \texttt{a5\_not\_near\_tight}]
$|A_5|=60$ and $\tfrac18\sqrt{60}<1$. Every $x\in A_5$ has $|C(x)|\ge1>\tfrac18\sqrt{60}$, so
\texttt{AttainedByA5Involution} is false. Every $x\ne1$ has $|C(x)|\ge2>2\cdot\tfrac18\sqrt{60}$.
\end{theorem}
\begin{proof}
$\sqrt{60}<8$ because $60<64$. Every subgroup has order $\ge1$. If $x\ne1$, then $1$ and $x$ are distinct
elements of $C(x)$.
\end{proof}

For the record (checked by brute force in the attached Python, not used in Lean), the centralizer orders in
$A_5$ are $3,4,5$ and $60$. All $15$ involutions have $|C(t)|=4$, so the involution ratio is
$4/\sqrt{60}\approx0.516$, far from $1/8$. Exact equality $k=\tfrac18\sqrt{|G|}$ would need
$|G|=64k^2$, and $60$ is not of that form.

\subsection*{Reading (R2): involutions in $\mathrm{PSL}(2,2^k)$}

Let $F$ be a finite field of characteristic $2$ with $q=|F|$. Recall
$\mathrm{PSL}(2,F)=\mathrm{SL}(2,F)/Z(\mathrm{SL}(2,F))$ (Mathlib's \texttt{Matrix.ProjectiveSpecialLinearGroup}).

\begin{lemma}[\texttt{eq\_one\_of\_sq\_eq\_one}, \texttt{center\_SL\_eq\_bot}, \texttt{toPSL}]
If $a\in F$ and $a^2=1$ then $a=1$. The center of $\mathrm{SL}(2,F)$ is trivial, so the quotient map
$\mathrm{SL}(2,F)\to\mathrm{PSL}(2,F)$ is an isomorphism.
\end{lemma}
\begin{proof}
$(a+1)^2=a^2+1=1+1=0$, so $a=-1=1$. By \texttt{SpecialLinearGroup.mem\_center\_iff} a central element is a
scalar $rI$ with $r^2=1$, hence $r=1$. A surjective homomorphism with trivial kernel is bijective.
\end{proof}

\begin{lemma}[\texttt{T\_sq}, \texttt{T\_ne\_one}, \texttt{entries\_of\_comm}]
Let $T=\begin{pmatrix}1&1\\0&1\end{pmatrix}\in\mathrm{SL}(2,F)$. Then $T^2=1\ne T$. If
$A=\begin{pmatrix}a&b\\c&d\end{pmatrix}\in\mathrm{SL}(2,F)$ commutes with $T$, then $a=d=1$ and $c=0$.
\end{lemma}
\begin{proof}
$T^2=\begin{pmatrix}1&2\\0&1\end{pmatrix}=1$ in characteristic $2$. $AT=\begin{pmatrix}a&a+b\\c&c+d\end{pmatrix}$
and $TA=\begin{pmatrix}a+c&b+d\\c&d\end{pmatrix}$, so $c=0$ and $a=d$; then $\det A=a^2=1$ gives $a=1$.
\end{proof}

\begin{lemma}[\texttt{card\_SL\_lower}, \texttt{card\_centralizer\_le}]
$|\mathrm{SL}(2,F)|\ge(q-1)q^2$, and the centralizer of $t=\mathrm{toPSL}(T)$ in $\mathrm{PSL}(2,F)$ has order
at most $q$.
\end{lemma}
\begin{proof}
$(a,b,c)\mapsto\begin{pmatrix}a&b\\c&a^{-1}(1+bc)\end{pmatrix}$ is an injection $F^\times\times F\times F\to
\mathrm{SL}(2,F)$. For the centralizer, $g\mapsto(\mathrm{toPSL}^{-1}g)_{01}$ is injective on $C(t)$: by the
previous lemma $\mathrm{toPSL}^{-1}g=\begin{pmatrix}1&b\\0&1\end{pmatrix}$ is determined by $b$. (In fact
$|C(t)|=q$ and $|G|=q(q^2-1)$; only the two inequalities are used.)
\end{proof}

\begin{lemma}[\texttt{psl\_facts}]
If $q\ge4$, then $G=\mathrm{PSL}(2,F)$ is simple, nonabelian and of even order, $t$ has order $2$,
$|C_G(t)|\le q$ and $|G|\ge(q-1)q^2$.
\end{lemma}
\begin{proof}
Simplicity is Mathlib's \texttt{Matrix.ProjectiveSpecialLinearGroup.rank\_two\_simple} (for $|F|\ge4$). The
order of $t$ is $2$ by \texttt{orderOf\_eq\_prime}, since $T^2=1\ne T$ and toPSL is injective; hence $2\mid|G|$.
If $G$ were abelian, then $C_G(t)=G$ and $(q-1)q^2\le|G|\le q$, which is impossible for $q\ge4$.
\end{proof}

\begin{theorem}[\texttt{not\_lowerBoundInvolutions}]
Reading (R2) with $c=1/8$ is false.
\end{theorem}
\begin{proof}
Take $F=\texttt{GaloisField 2 7}$, so $q=128$ (\texttt{GaloisField.card}), and $x=t$. Apply
\texttt{ratio\_aux} with $q=128$, $r=127$, $c=\tfrac18$: $c^2r=127/64>1$, so
$|C_G(t)|<\tfrac18\sqrt{|G|}$. Numerically $|G|=128\cdot16383=2097024$ and $\tfrac18\sqrt{|G|}\approx181>128$.
\end{proof}

\begin{theorem}[\texttt{psl\_centralizer\_ratio\_small}]
For every real $c>0$ there is $k$ such that $G=\mathrm{PSL}(2,\texttt{GaloisField 2 k})$ is simple, nonabelian
and of even order and has an involution $t$ with $|C_G(t)|<c\sqrt{|G|}$.
\end{theorem}
\begin{proof}
Let $k=\lceil1/c^2\rceil+2$ and $q=2^k\ge4$. Since $q-1\ge k>1/c^2$, apply \texttt{ratio\_aux} with $r=q-1$.
\end{proof}

\begin{theorem}[\texttt{conjecture\_2291\_false}, main theorem]
The conjunction of: $\neg$\texttt{LowerBoundAllElements}; the every-$c>0$ statement for (R1);
$\neg$\texttt{LowerBoundInvolutions}; the every-$c>0$ statement for (R2); $\neg$\texttt{AttainedByA5Involution};
and the factor-$2$ slack at every nonidentity element of $A_5$.
\end{theorem}

\section{Lean formalization and axiom audit}

File \texttt{lean/Conjecture2291/Basic.lean} (namespace \texttt{C2291}, Lean 4 / Mathlib v4.33.1, only
\texttt{import Mathlib}). Theorems: \texttt{finRotate\_mem\_alternatingGroup},
\texttt{card\_centralizer\_finRotate\_perm}, \texttt{card\_centralizer\_cyc\_le}, \texttt{cyc\_ne\_one},
\texttt{two\_mul\_card\_A'}, \texttt{factorial\_lower}, \texttt{A\_hyps}, \texttt{ratio\_aux},
\texttt{alternating\_centralizer\_ratio\_small}, \texttt{not\_lowerBoundAllElements}, \texttt{card\_A5},
\texttt{sqrt60\_div8\_lt\_one}, \texttt{a5\_not\_near\_tight}, \texttt{not\_attainedByA5Involution},
\texttt{eq\_one\_of\_sq\_eq\_one}, \texttt{center\_SL\_eq\_bot}, \texttt{toPSL}, \texttt{T\_sq}, \texttt{T\_ne\_one},
\texttt{entries\_of\_comm}, \texttt{card\_SL\_lower}, \texttt{card\_centralizer\_le}, \texttt{psl\_facts},
\texttt{not\_lowerBoundInvolutions}, \texttt{psl\_centralizer\_ratio\_small}, and the main theorem \texttt{conjecture\_2291\_false}. \texttt{lake build} succeeds, and
\texttt{\#print axioms C2291.conjecture\_2291\_false} reports
\texttt{[propext, Classical.choice, Quot.sound]}. There is no \texttt{sorry} and no \texttt{native\_decide}.

\end{document}
